\section{Discussion}
\label{sec:discussion}

% The error analysis in Section~\ref{sec:error-analysis} reveals that 14 of 21 execution-stage failures stem from two tractable root causes: invisible pybind11 type signatures in API documentation and insufficient assertion-level validation feedback. These findings directly inform the design of more robust evaluation harnesses and type-aware retrieval strategies.
% % 第\ref{sec:error-analysis}节的错误分析表明，21个执行阶段失败中有14个源于两个可处理的根本原因：API文档中不可见的pybind11类型签名，以及断言级验证反馈不足。这些发现直接为更健壮的评估框架和类型感知检索策略的设计提供了依据。
The error analysis in Section~\ref{sec:error-analysis} shows that the 21 execution-stage failures mainly arise from API misuse and incorrect API composition. These failures are largely caused by undocumented API contracts, missing usage-level knowledge, and insufficient validation feedback, motivating type-aware retrieval and more informative execution feedback.
% 第\ref{sec:error-analysis}节的错误分析表明，21个执行阶段失败主要源于API误用和错误的API组合。这些失败大多由未文档化的API契约、用法级知识的缺失以及验证反馈不足引起，从而说明了类型感知检索和更充分执行反馈的必要性。

Beyond these specific improvements, our results carry broader implications. The dominance of sandbox execution (+[TBD] pp) over static knowledge suggests that for domains with long-tailed, under-docu\-mented APIs, execution grounding is not merely beneficial but essential. Our theoretical characterization (Section~\ref{sec:grounding-theory}) elevates this from an empirical observation to a structural principle: the grounding gap is bounded by observation fidelity, so execution channels—not larger priors—are the binding constraint in such domains. This paradigm may generalize to other domain-specific scripting tasks with similar characteristics.
% 除了这些具体改进之外，我们的结果具有更广泛的意义。沙盒执行（+[TBD]个百分点）相对于静态知识的主导地位表明，对于具有长尾、文档不完善API的领域，执行驱动不仅是锦上添花，而是关键所在。这一范式可能推广到其他具有相似特征的领域特定脚本任务。

Several limitations also warrant acknowledgment. First, performance varies significantly across base LLMs (16.7\%–84.8\%), indicating that the underlying model remains a decisive factor. Second, the interactive bridge is single-threaded, limiting concurrent deployment. Third, while realistic, the primary benchmark is limited to the 158-task PyAether slice; the remaining SKILL and Tcl slices of EDA-Scripting-Bench are under construction, and every reported figure reflects the PyAether-only slice and will be regenerated on the full benchmark. Fourth, although each configuration was evaluated over five independent runs, the benchmark size limits statistical power to certify small effects.
% 几个局限性也值得说明。第一，不同基座LLM表现差异显著（16.7%–84.8%），表明底层模型仍然是决定性因素。第二，交互桥接是单线程的，限制了并发部署。第三，主评测集虽然真实，但仅限于158个PyAether切片；EDA-Scripting-Bench的SKILL与Tcl切片正在建设中，所有报告数据均为PyAether切片结果，将在完整评测集上重新生成。第四，尽管每个配置运行了五次独立实验，但评测集规模限制了检验小效应的统计功效。